\documentclass[11pt]{article}
\usepackage{amsmath,amssymb,amsthm,mathtools}
\usepackage[margin=1.05in]{geometry}
\usepackage[hidelinks]{hyperref}
\newtheorem{theorem}{Theorem}[section]
\newtheorem{proposition}[theorem]{Proposition}
\newtheorem{lemma}[theorem]{Lemma}
\newtheorem{corollary}[theorem]{Corollary}
\newcommand{\Pp}{\mathbb P}
\newcommand{\E}{\mathbb E}
\newcommand{\SSer}{\mathfrak S}
\newcommand{\ind}{\mathbf 1}
\newcommand{\Prob}{\mathbb P_k}
\newcommand{\dd}{\,\mathrm d}
\newcommand{\doilink}[1]{\href{https://doi.org/#1}{doi:\,#1}}
\title{Extremes of Goldbach counts on progressions\\ with a prescribed endpoint}
% Enter the author's approved public byline here before arXiv submission.
\author{}
\date{}
\begin{document}
\maketitle
\begin{abstract}
Fix a large even integer \(N\). We study the number \(R(N-ih)\) of ordered
representations as a sum of two primes along a progression leading up to \(N\).
The step \(h\) is a multiple of the product of the primes up to
\(w=\log\log N\). If \(1\le i\le(\log N)^D\), the maximum of
\((\log w)\log(R(N-ih)/X_N)\) converges in probability to \(\log(1+D)\),
where \(X_N\) is the Hardy--Littlewood main term with its local factors
truncated at \(w\). At each fixed level \(0<t<\log(1+D)\), the number
of exceedances is \((\log N)^{D+1-e^t+o_{\mathbb P}(1)}\).
Both statements hold when the multiplier is sampled uniformly from the
eligible integers or the eligible primes, uniformly over prescribed
endpoints and step intervals \((H,2H]\) with
\(N^{8/33+\varepsilon}\le H\le N^{1-\varepsilon}\).
The proof transfers a large-deviation estimate for reciprocal prime divisors
to actual representation counts using the averaged Goldbach theorem of
Matom\"aki, Radziwi{\l}{\l} and Tao.
\end{abstract}
\section{The endpoint and the sampling law}

Let \(R(n)\) be the number of ordered pairs of primes with sum \(n\). The
Hardy--Littlewood main term for \(R(n)\), for even \(n\), is
\(n\SSer(n)/(\log n)^2\), where
\begin{equation}\label{eq:series}
 C_2=\prod_{p>2}\left(1-\frac{1}{(p-1)^2}\right),\qquad
 \SSer(n)=2C_2\prod_{\substack{p\mid n\\p>2}}\frac{p-1}{p-2}.
\end{equation}
Along a progression with step divisible by all small primes,
the corresponding factors in \eqref{eq:series} stay fixed. The larger
prime divisors depend on the position in the progression. We determine
how large their combined effect becomes among a growing number of actual
Goldbach counts.

All logarithms are natural. Let \(N\) tend to infinity through the even
integers. Write
\begin{equation}\label{eq:parameters}
 L=\log N,\quad w=\log L,\quad W=\prod_{p\le w}p,\quad
 K=\lfloor L^D\rfloor,\quad
 X_N=\frac{N}{L^2}\,2C_2\prod_{\substack{2<p\le w\\p\mid N}}
                                 \frac{p-1}{p-2},
\end{equation}
where \(D>0\) is fixed. Set \(\sigma=8/33\). Fix
\(0<\varepsilon<(1-\sigma)/2\) and choose
\(N^{\sigma+\varepsilon}\le H\le N^{1-\varepsilon}\). Put \(Y=H/W\).
Choose \(k\) uniformly from the integers in \((Y,2Y]\), or uniformly
from the primes in that interval. In both cases let \(h=Wk\). The
probability and expectation below refer solely to this choice of \(k\).
In particular, all estimates are uniform over the even endpoints \(N\)
and admissible \(H\). For \(1\le i\le K\), put
\[
 n_i=N-ih,\qquad Z_i=(\log w)\log\frac{R(n_i)}{X_N}.
\]
Set \(Z_i=-\infty\) if \(R(n_i)=0\).

\begin{theorem}\label{thm:main}
For either sampling law,
\[
  \min_{1\le i\le K}Z_i \ \xrightarrow{\ \mathrm{P}\ }\ 0,
  \qquad
  \max_{1\le i\le K}Z_i \ \xrightarrow{\ \mathrm{P}\ }\ \log(1+D).
\]
For \(0<t<\log(1+D)\), let \(C_t=\#\{1\le i\le K:Z_i\ge t\}\).
Then
\begin{equation}\label{eq:profile}
        C_t=L^{D-(e^t-1)+o_{\mathrm P}(1)} .
\end{equation}
Precisely, for every fixed \(\eta>0\), the probability that
\(L^{D+1-e^t-\eta}\le C_t\le L^{D+1-e^t+\eta}\) tends to one.
\end{theorem}

For example, take \(K=\lfloor\log N\rfloor\). The largest value of
\(Z_i\) then tends in probability to \(\log 2\), and
\(L^{1/2+o_{\mathrm P}(1)}\) positions reach the level \(\log(3/2)\).
The endpoint \(N\) may be chosen before the step. The theorem concerns the
positive indices \(i\ge1\) and leaves \(R(N)\) undetermined.

Erd\H{o}s and K\'atai studied maxima of a nonnegative strongly additive
function along consecutive integers with the starting integer varying.
Their Theorem~4 treats the prime values \(g(p)=1/p\)
\cite{EK}. Weingartner estimated tails of the limiting distributions
of \(\sigma(n)/n\) and \(n/\varphi(n)\) by analysing high moments of
Euler products \cite{Wei}. Languasco computed fixed-order moments of
the full binary Goldbach singular series over \(n\le x\) and applied them
to Goldbach numbers in short intervals \cite{Lan}.
Here the cutoff \(w\) varies with the endpoint,
the step supplies the randomness, and all positions share this step. The
proof retains the factors at primes dividing both the index \(i\) and
the prescribed endpoint \(N\). These factors can be large for an
individual \(i\). Averaging them over the index range preserves the
exponential rate.

\section{Transferring the averaged Goldbach theorem}

Let \(\Lambda\) be the von Mangoldt function and write
\(G(n)=\sum_{p+q=n}\log p\log q\), the sum running over ordered prime
pairs. We use Theorem~1.3(iv) of Matom\"aki, Radziwi{\l}{\l} and Tao
\cite{MRT}. For fixed \(\varepsilon'>0\) in the range of that theorem,
\(x^{\sigma+\varepsilon'}\le Z\le x^{1-\varepsilon'}\), and any
\(E>0\), one has
\begin{equation}\label{eq:mrt}
  \sum_{a+b=n}\Lambda(a)\Lambda(b)
     = n\SSer(n)+O_{E,\varepsilon'}(x/\log^E x)
\end{equation}
for all but \(O_{E,\varepsilon'}(Z/\log^E x)\) even integers in
\([x,x+Z]\). This is the analytic input about binary Goldbach counts.

\begin{lemma}\label{lem:transfer}
Let \(y=L^2\), and put
\[
        T_i(k)=\sum_{\substack{w<p\le y\\p\mid n_i}}\frac1p .
\]
For each fixed \(A>0\), outside a set of \(k\) of probability
\(O_{A,D,\varepsilon}(L^{-A})\) under either law, all \(R(n_i)>0\) and
\begin{equation}\label{eq:transfer}
        \max_{i\le K}|Z_i-(\log w)T_i(k)|=o(1).
\end{equation}
\end{lemma}

\begin{proof}
For a fixed \(i\le K\), apply \eqref{eq:mrt} with
\(x=N-2iH\), \(Z=iH\) and \(\varepsilon'=\varepsilon/2\).
Here \(x\sim N\), and the factor \(i\le L^D\) fits inside the slack
at the upper end of the permissible range. Each exceptional \(n_i\)
determines at most one \(k\). Summing over \(i\) costs
\(O_{E,D,\varepsilon}(K^2H/L^E)\) choices. The prime number theorem
gives \(\log W=(1+o(1))w\). The two sampling sets have sizes
\(\asymp Y\) and \(\asymp Y/\log Y\), respectively. Choosing \(E\)
large compared with \(A\) and \(D\) proves the probability estimate.

There are \(O(\sqrt N L)\) nontrivial prime powers up to \(N\).
Their contribution for a given \(n_i\) is \(O(\sqrt N L^3)\).
For \(p,q>N/L^4\), the weight \((\log p)(\log q)\) equals
\(L^2(1+O(w/L))\). At most \(2N/L^4+O(1)\) ordered pairs have a smaller
endpoint. Since \(\SSer(n_i)\ge2C_2\), on the retained steps these estimates
give, uniformly in \(i\),
\[
 R(n_i)=\frac{N\SSer(n_i)}{L^2}
             \left(1+O\left(\frac{w}{L}+\frac{KH}{N}\right)\right).
\]
For \(p\le w\), membership in the prime divisors of \(n_i\) is fixed by
\(N\). Above \(w\), replacing
\(\log((p-1)/(p-2))\) by \(1/p\) costs \(O(\sum_{m>w}m^{-2})=O(1/w)\).
Finally \(\sum_{p>y,\,p\mid n_i}1/p\le L/(y\log y)\). Multiplication
by \(\log w\) sends all these errors to zero.
\end{proof}

\section{Large values from primes above the cutoff}

Fix \(\lambda>1\), and put \(\tau=\lambda w\log w\) and
\(b_p=\exp(\tau/p)-1\). The following product measures the exponential
cost of forcing prime divisors just above \(w\).
\begin{lemma}\label{lem:product}
For \(y=L^2\),
\[
 A_\lambda=\prod_{w<p\le y}\left(1+\frac{b_p}{p}\right),\qquad
 \log A_\lambda=w(\lambda\log\lambda-\lambda+1)+o_\lambda(w).
\]
\end{lemma}
\begin{proof}
Fix \(0<\delta<(\lambda-1)/2\). For \(p=uw\) and
\(1\le u\le\lambda-\delta\),
\[
 \log(1+b_p/p)=(\lambda/u-1)\log w+O_\lambda(1)
\]
uniformly as \(w\to\infty\). The prime number theorem and partial summation
give a contribution
\(w\int_1^{\lambda-\delta}(\lambda/u-1)\dd u+o_{\lambda,\delta}(w)\).
For \((\lambda-\delta)w<p\le(\lambda+\delta)w\), each logarithm is
\(O_\lambda(\log w)\). The prime number theorem bounds their sum by
\(O_\lambda(\delta w)+o_{\lambda,\delta}(w)\). For \(p>(\lambda+\delta)w\),
use \(e^v-1\le ve^v\) for \(v\ge0\). The remaining sum is at most
\[
 \tau w^{\lambda/(\lambda+\delta)}
          \sum_{m>(\lambda+\delta)w}m^{-2}=o_{\lambda,\delta}(w).
\]
Letting \(\delta\) decrease to zero evaluates the integral as stated.
\end{proof}

\begin{lemma}\label{lem:moment}
For either sampling law and each fixed \(\lambda>1\),
\[
 \sum_{i\le K}\E_k e^{\tau T_i(k)}
       \le K A_\lambda e^{o_\lambda(w)},\qquad
 \E_k e^{\tau T_1(k)}\le A_\lambda e^{o_\lambda(w)}.
\]
\end{lemma}
\begin{proof}
When \(p\mid i\), the condition \(p\mid N-iWk\) holds for every \(k\)
if \(p\mid N\), and for no \(k\) otherwise. Every other prime in \((w,y]\)
specifies a single residue class for \(k\bmod p\). In the product
\(\prod_{w<p\le y}(1+b_p\ind_{p\mid n_i})\), take out the constant
factors at \(p\mid(i,N)\) and expand over the primes \(p\nmid i\).
For a squarefree product \(d\), all its congruences specify one class
modulo \(d\) by the Chinese remainder theorem. The sum of all expansion
weights, including the factors forced by \(p\mid(i,N)\), is at most
\[
 \prod_{w<p\le y}(1+b_p)
    =\exp\left(\tau\sum_{w<p\le y}1/p\right)
    \le \exp(O_\lambda(w^2\log w))=Y^{o(1)}.
\]
For integer \(k\), a class contains \(Y/d+O(1)\) points in the interval.
Thus its expected contribution uses \(1/d\), with total additive error
\(Y^{-1+o(1)}\) after normalization by \(Y+O(1)\).

For prime \(k\), take \(d\le\sqrt Y\) first. A non-reduced class
contains no eligible prime, since every prime factor of \(d\) is at most
\(y<Y\). For a reduced class the Brun--Titchmarsh inequality
\cite[Theorem~2, (1.10)]{MV} bounds its probability by
\(O(1/\varphi(d))\), using the prime number theorem for the denominator
and treating \(d=1\) directly. Uniformly over these \(d\),
\(d/\varphi(d)\ll\log y=e^{o(w)}\) by Mertens' theorem. For larger \(d\)
the integer count bounds the probability by
\(O((Y^{-1/2}+Y^{-1})\log Y)\). With the expansion weights, its total
contribution is \(Y^{-1/2+o(1)}\).

Combining the two estimates and retaining all the positive local factors
in \(A_\lambda\), we obtain, under either law,
\[
 \E_k e^{\tau T_i(k)}
   \le e^{o_\lambda(w)} A_\lambda
        \prod_{\substack{w<p\le y\\p\mid(i,N)}}g_p+o(1),
 \qquad g_p=\frac{1+b_p}{1+b_p/p}\le\min(p,1+b_p).
\]
The errors here are uniform in \(i\) and \(N\). To sum the product over
\(i\le K\), expand it in the nonnegative factors \(g_p-1\). A squarefree
product \(d\) divides \(\lfloor K/d\rfloor\) of these indices. Consequently
the average is at most
\[
 \prod_{\substack{w<p\le y\\p\mid N}}\left(1+\frac{g_p-1}{p}\right)
        =e^{o_\lambda(w)} .
\]
Indeed, the factors with \(p\le2\lambda w\) are at most \(2\), and there
are \(O_\lambda(w/\log w)\) of them. The sum of the remaining logarithms
is at most \(\sum_{p>2\lambda w} b_p/p\ll_\lambda \sqrt w\log w\),
by the same exponential bound as in Lemma~\ref{lem:product}. For \(i=1\)
the product over common divisors is empty.
\end{proof}

\begin{proof}[Proof of Theorem~\ref{thm:main}]
Write \(U_i=(\log w)T_i(k)\). For a fixed \(s>0\) take
\(\lambda=e^s\) in Lemmas~\ref{lem:product} and~\ref{lem:moment}.
Markov's inequality yields
\begin{align}
 \E_k\#\{i\le K:U_i\ge s\}
   &\le L^{D+1-e^s+o(1)}, \label{eq:upper}\\
 \Prob(U_1\ge s)&\le L^{1-e^s+o(1)}. \nonumber
\end{align}
Here the exponential penalty is \(\exp(-\lambda sw)\).
Since \(U_i\ge0\), the second inequality gives the limit for the minimum.
For \(s>\log(1+D)\), the first gives the upper bound for the maximum.

For \(1<c<1+D\), let \(P_c=\prod_{w<p\le cw}p\). The prime number theorem
and partial summation give
\[
 \log P_c=(c-1+o(1))w,\qquad
 (\log w)\sum_{w<p\le cw}1/p=\log c+o(1).
\]
The first estimate implies \(P_c=o(K)\). For integer multipliers the
probability of a common factor of \(k\) and \(P_c\) is at most
\(\sum_{w<p\le cw}(1/p+O(1/Y))=O_c(1/\log w)\). Every prime multiplier
in the sampling interval exceeds \(cw\). Whenever \((k,P_c)=1\), the class
\[
             i\equiv N(Wk)^{-1}\pmod{P_c}
\]
contains \(K/P_c+O(1)=L^{D+1-c+o(1)}\) indices in \([1,K]\).
Every prime dividing \(P_c\) divides \(n_i\) at each such index, so
\(U_i\ge\log c+o(1)\). Taking \(c\) arbitrarily close to \(1+D\) from
below gives the lower bound for the maximum.

For a level \(0<t<\log(1+D)\), the same construction with \(c>e^t\)
supplies the lower bound in \eqref{eq:profile} by taking \(c\) as close to
\(e^t\) as the prescribed exponent tolerance requires. Its strict margin
\(\log c-t\) absorbs the error in Lemma~\ref{lem:transfer}. For the upper
bound apply \eqref{eq:upper} at \(s<t\) sufficiently close to \(t\),
then apply Markov's inequality to the count. Lemma~\ref{lem:transfer}
bounds the difference between all the \(U_i\) and \(Z_i\) simultaneously
on a set of probability tending to one. This also transfers the two
extreme-value limits to the representation counts.
\end{proof}

\section{Steps with a common singular series}\label{sec:sparse}

Write \(S_w(N)=L^2X_N/N\) for the truncated singular series in
\eqref{eq:parameters}. We can select integer multipliers for which a
growing initial segment of the progression has this same factor to
arbitrary inverse-logarithmic accuracy.

\begin{theorem}\label{thm:sparse}
Fix \(A,c>0\) and set \(r=\lfloor cw/\log w\rfloor\).
For every sufficiently large even \(N\) and every admissible \(H\),
at least \(Y L^{-c-o(1)}\) integers \(k\in(Y,2Y]\) satisfy
\begin{equation}\label{eq:sparse}
    G(n_i)=n_i S_w(N)+O_{A,c,\varepsilon}(N/L^A)
    \qquad (1\le i\le r).
\end{equation}
\end{theorem}

\begin{proof}
Put \(z=L^{A+4}\). For each prime \(w<p\le z\), exclude the residues
of \(k\bmod p\) for which \(p\mid N-iWk\) for at least one
\(1\le i\le r\). For large \(N\) we have \(r<w/2\). Thus \(iW\) is
invertible modulo \(p\), and the number of excluded residues is
\[
 \rho(p)=
 \begin{cases}
  1,&p\mid N,\\
  r,&p\nmid N .
 \end{cases}
\]
Indeed, the residues \(N(iW)^{-1}\bmod p\) coincide at zero in the
first case and are distinct in the second. Let
\[
 V=\prod_{w<p\le z}(1-\rho(p)/p),\qquad
 U=\sum_{w<p\le z}\rho(p)/p .
\]
By Mertens' estimate,
\(U\le r(\log\log z-\log\log w+O(1))\le(c+o(1))w\).
Also \(\sum_{w<p\le z}(\rho(p)/p)^2\ll r^2/w=o(w)\). Expanding the
logarithm gives \(\log V\ge-(c+o(1))w\).

We next count the surviving integers in \((Y,2Y]\). Fix \(B>c+1\).
For a squarefree product \(d\) of the sieving primes, put
\(\rho(d)=\prod_{p\mid d}\rho(p)\). The Chinese remainder theorem gives
\(Y\rho(d)/d+O(\rho(d))\) multipliers in the intersection of the
corresponding excluded classes. Truncate inclusion--exclusion at the
consecutive orders \(t-1,t\), with \(t\) odd. The two Bonferroni bounds
for the model proportion bracket \(V\) and differ by at most
\(U^t/t!\). Choose \(t=Cw+O(1)\) with \(C\) sufficiently large in
terms of \(A,c,B\). Then \(U^t/t!\le(eU/t)^t\ll L^{-B}\).
For either truncation the total rounding error, divided by \(Y\), is
at most a constant times
\[
 Y^{-1}\!\!\sum_{\substack{d\\\omega(d)\le t}}\rho(d)
 \le Y^{-1}(1+rz)^t
 =Y^{-1}\exp(O_{A,c,B}(w^2))=o(L^{-B}),
\]
where the sum runs over the squarefree products just specified and
\(\log Y\gg_\varepsilon L\). The integer count is therefore
\(Y(V+O_{A,c,B,\varepsilon}(L^{-B}))\).

For a surviving \(k\), every prime divisor of \(n_i\) up to \(z\)
already lies below or at \(w\), where \(n_i\equiv N\pmod W\) fixes
the local factors. For all \(i\le r\) we have
\[
 0\le\log\frac{\SSer(n_i)}{S_w(N)}
   \ll\sum_{\substack{p>z\\p\mid n_i}}\frac1p
   \le \frac{L}{z\log z}.
\]
Mertens' product estimate gives \(S_w(N)\ll\log w\).
Thus \(n_i\SSer(n_i)=n_i S_w(N)+O_A(N/L^A)\).

Finally apply \eqref{eq:mrt} on \([N-2iH,N-iH]\) for each \(i\le r\).
As in Lemma~\ref{lem:transfer}, the exceptional multipliers number
\(O_{E,\varepsilon}(r^2H/L^E)\). Since \(W=L^{1+o(1)}\) and \(H=WY\),
we may choose \(E>A+B+3\) to bound this number by \(O(YL^{-B})\).
Removing higher prime powers costs \(O(\sqrt N L^3)\) at each target.
The resulting count of multipliers is still at least \(Y L^{-c-o(1)}\),
and each satisfies \eqref{eq:sparse}.
\end{proof}

When \(N\) runs through powers of two, each sieving prime excludes exactly
\(r\) classes. In this case \(U=(c+o(1))w\) and the sieve itself retains
\(Y L^{-c+o(1)}\) multipliers. This measures the size of the particular
selection used in the proof.

\section*{Acknowledgements and availability}

OpenAI Codex assisted in developing the arguments, preparing the manuscript
and writing the companion computational checks. The
\href{https://github.com/deltavra/goldbach-representations}{companion repository}
includes finite arithmetic examples and a separate Lean certificate for
Goldbach partitions of even integers through \(2^{32}\). The
asymptotic results in this paper are given by the arguments above.
This manuscript is distributed under the
\href{https://creativecommons.org/licenses/by/4.0/}{Creative Commons
Attribution 4.0 International license}.

\begin{thebibliography}{9}

\bibitem{EK}
P.~Erd\H{o}s and I.~K\'atai,
On the maximal value of additive functions in short intervals and on some
related questions,
\emph{Acta Math. Acad. Sci. Hungar.} \textbf{35} (1980), 257--278.
\doilink{10.1007/BF01896843}.
Correction, \emph{ibid.} \textbf{37} (1981), 499.
\doilink{10.1007/BF01895152}.

\bibitem{Lan}
A.~Languasco, A singular series average and Goldbach numbers in short
intervals, \emph{Acta Arith.} \textbf{83} (1998), 171--179.
\doilink{10.4064/aa-83-2-171-179}.

\bibitem{MRT}
K.~Matom\"aki, M.~Radziwi{\l}{\l} and T.~Tao,
Correlations of the von Mangoldt and higher divisor functions I. Long shift
ranges, \emph{Proc. Lond. Math. Soc.} (3) \textbf{118} (2019), 284--350.
\doilink{10.1112/plms.12181}.

\bibitem{MV}
H.~L. Montgomery and R.~C. Vaughan, The large sieve,
\emph{Mathematika} \textbf{20} (1973), 119--134.
\doilink{10.1112/S0025579300004708}.

\bibitem{Wei}
A.~Weingartner, The distribution functions of \(\sigma(n)/n\) and
\(n/\varphi(n)\), \emph{Proc. Amer. Math. Soc.} \textbf{135} (2007),
2677--2681. \doilink{10.1090/S0002-9939-07-08771-0}.

\end{thebibliography}

\end{document}
